\section{Executive Investment Recommendation}

Based on the portfolio optimization, historical performance, downside-risk measures, and stress-testing results, the Maximum Sharpe Ratio portfolio is recommended. The portfolio allocates 55.44 percent to NVIDIA (NVDA), 26.14 percent to Bitcoin (BTC), and 18.42 percent to Manila Electric Company (MER), delivering an expected annual return of 58.10 percent, an annualized volatility of 38.25 percent, and a Sharpe ratio of 1.3815 at the 5.25 percent risk-free rate, the highest risk-adjusted profile in the feasible universe. Figure~\ref{fig:portfolio_allocation} contrasts this allocation with the GMV portfolio, and Figure~\ref{fig:efficient_frontier} locates it at the tangency point of the efficient frontier.

NVIDIA contributes the strongest sample-period expected return (71.65 percent annualized) and, together with Bitcoin (61.49 percent annualized), drives the portfolio's return profile. Bitcoin's inclusion rests on its high expected return and its low correlation with the Philippine constituents (MER--BTC 0.0286), while NVIDIA pairs with the domestic block at near-zero correlation (MER--NVDA 0.0121, MER--META 0.0172) and at moderate co-movement with Bitcoin (NVDA--BTC 0.3204). Manila Electric Company anchors the allocation as the sole domestic constituent, combining a positive expected return (12.49 percent annualized) with near-zero correlation to both growth assets. TEL, AEV, and META receive zero weight in the long-only tangency solution.

Downside risk is governed by the empirical tail measures reported in Table~\ref{tab:var_summary}. The recommended portfolio's daily 95 percent Historical VaR is -3.39 percent and its 95 percent Historical Expected Shortfall is -5.11 percent; at the 99 percent level these reach -5.64 percent and -8.55 percent. The empirical 99 percent Expected Shortfall (8.55 percent) exceeds the Gaussian parametric estimate (6.19 percent) by 236 basis points per day, so historical measures, rather than the normal approximation, should set the risk budget.

Stress testing confirms the trade-off between return and capital preservation. Under the five scenarios in Figure~\ref{fig:stress_test} and Table~\ref{tab:stress_summary}, the Maximum Sharpe portfolio loses 28.83 percent in the 2020 COVID crash, 24.58 percent in the technology-sector selloff, 18.65 percent in the 2022 rate-hike shock, and 14.40 percent in the 2022 crypto winter, with a -45.30 percent loss under the hypothetical 1.5x Black Swan scenario. The GMV portfolio contains the same shocks at materially smaller magnitudes (-22.55 percent COVID, -35.89 percent Black Swan), illustrating its role as the defensive alternative.

Disciplined monthly rebalancing is required to manage the portfolio's concentration: NVDA and BTC together represent 81.58 percent of the recommended allocation, and un-rebalanced equal-weight exposure drifted to 61.75 percent NVDA and 24.77 percent BTC by December 2025. Over the sample, monthly rebalancing reduced volatility from 35.12 percent to 24.18 percent and the peak drawdown from -58.21 percent to -43.20 percent at the cost of cumulative return (384.85 percent versus 743.04 percent). For investors prioritizing capital preservation, the GMV portfolio allocates TEL 31.01 percent, MER 31.23 percent, AEV 16.85 percent, NVDA 3.15 percent, META 12.86 percent, and BTC 4.90 percent, delivering a 16.03 percent expected return at 20.22 percent volatility (Sharpe 0.5332). Portfolio choice ultimately reflects the investor's willingness and capacity to bear losses: the Maximum Sharpe portfolio suits investors seeking higher returns despite greater risk, while the GMV portfolio suits investors who cannot tolerate peak stress losses near -45 percent.
